% !TITLE 己亥杂诗·其五
% !AUTHOR 龚自珍
% !DYNASTY 近代
% !GENRE 七言绝句
% !ORDER 1

\begin{poem}
浩荡\textsuperscript{1}离愁白日斜，吟鞭\textsuperscript{2}东指即天涯\textsuperscript{3}。
落红\textsuperscript{4}不是无情物，化作春泥更护花\textsuperscript{5}。
\end{poem}

\begin{annotations}
\notetext{1}{浩荡：广阔浩大。形容离愁不是一点半点，而是弥漫天地。}
\notetext{2}{吟鞭：诗人的马鞭。一边吟诗一边挥鞭策马东行。}
\notetext{3}{即天涯：就是天涯。离开京城（北京）东归杭州，感觉像去往天涯一般遥远。}
\notetext{4}{落红：落花。凋谢的花瓣。诗人以落花自比——离开官场如同花儿凋谢。}
\notetext{5}{化作春泥更护花：落花融入泥土变成春天的肥料，反而能养护新开的花。比喻自己虽然离开官场，仍会以另一种方式报效国家。}
\end{annotations}

\begin{appreciation}
《己亥杂诗》是龚自珍四十八岁那年辞官南归，一路写下的组诗，一共三百一十五首，这一首是其中流传最广的几首之一。龚自珍是公认的近代开风气的人，旧时代已经散架，他最早闻到了味道。

“浩荡离愁白日斜”。浩荡本是形容大水的气势，用在这里，离愁大得像要漫过堤岸。白日斜——太阳已经西斜，人还得赶路。这不是普通的离别：辞官，等于把二十年的仕途连根拔起，四十八岁的人，前路只剩一个“归”字。“吟鞭东指即天涯”——策马向东回杭州，一边走一边吟诗，这“天涯”不在天边，就在马鞭指着的方向：离开了京城，从此就是另一种人生。

前两句是沉甸甸的愁，后两句却忽然翻出一层亮色：“落红不是无情物，化作春泥更护花”。落红就是落花。古人写落花，多半是伤春的套话：花谢了，春天走了，人也老了。龚自珍偏不认这个账。他说落花不是无情之物——它落进土里，化作春泥，来年再滋养新开的花。他拿自己比作那片落红：官场是回不去了，还可以回杭州讲学、著书，把一肚子学问交给年轻人。凋谢不是结束，是换了种活法。

这一句的分量，放进整本书里才看得清。屈原在《离骚》里佩兰纫芷，以香草装点自己，那是中国文学香草传统的源头：美好的人与美好的草木互相印证。龚自珍不佩香草了，他把自己化成了土——同样是芬芳，却让给了下一季。至于《生年不满百》里那桩心事，“生年不满百，常怀千岁忧”，生命太短，忧虑太长，那时的诗人回答“何不秉烛游”；龚自珍却有另一种解法：人虽只有一世，花却可以开两季——自己谢了，春泥里还存着下一春。

你以后也会有很多告别：毕业、搬家、离开要好的朋友。舍不得是应当的，但别把告别当成结束。落红不是无情物——这句话写于一百多年前，说的也是你将来要说的事。
\end{appreciation}
